\section{Ingeniería inversa de la autenticación MD5 de RIP}\label{sec:reveng}
Cada captura es una Respuesta RIPv2 enviada a la dirección multicast 224.0.0.9 sobre UDP/520. La autenticación de tipo~3 coloca una cabecera de auth como primera entrada de la tabla de rutas y añade un tráiler de auth que transporta el digest (16 bytes). La captura se descompone como en la \autoref{tab:layout}, ese digest de 16 bytes es el campo que el ataque debe reproducir.

\begin{table}[H]
\centering
\footnotesize
\begin{tabularx}{\linewidth}{l l X}
\hline
\textbf{Región} & \textbf{Bytes} & \textbf{Contenido (captura tamaño~06)}\\
\hline
Cabecera RIP        & 4  & command=2, version=2, routing-domain=0\\
Cabecera auth (RTE) & 20 & AFI=0xFFFF, auth-type=3, pkt-len=44, key-id=1, auth-len=20, seq, reservado\\
Entrada de ruta     & 20 & la ruta anunciada\\
Tráiler de auth     & 20 & AFI=0xFFFF, type=0x0001, seguido del digest de 16 bytes\\
\hline
\end{tabularx}
\caption{Formato del paquete autenticado RIPv2/RFC~2082 (payload de 64 bytes).}
\label{tab:layout}
\end{table}

\paragraph{Construcción del digest (hallazgo principal).} El digest capturado corresponde a keyed-MD5 de RFC~2082. El campo de digest de 16 bytes se sustituye primero por la clave y se calcula un MD5 simple sobre todo el paquete. Sea $D$ la región desde la cabecera RIP hasta los campos \texttt{AFI}/\texttt{type} del tráiler (todo salvo los 16 bytes del digest) y $K_{16}$ la clave ASCII rellenada con bytes 0-16, como en la ecuación~\eqref{eq:keyed}.
\begin{equation}
\text{digest} \;=\; \mathrm{MD5}\bigl(\, D \,\|\, K_{16} \,\bigr).
\label{eq:keyed}
\end{equation}
Se probaron ambas construcciones (keyed-MD5 y HMAC-MD5), solo la ecuación~\eqref{eq:keyed} sobre un alfabeto base-36 (letras minúsculas y dígitos, no hexadecimal) reprodujo el digest capturado, el espacio hexadecimal se agotó sin coincidencia (keyed-MD5 en 7.0~s, HMAC-MD5 en 20.4~s). Por lo tanto, base-36 con keyed-MD5 es la construcción confirmada.

\paragraph{Claves recuperadas y validación.} Se recuperó una clave por captura, pero ninguna se aceptó por coincidir con un solo paquete. Se recalculó el digest keyed-MD5 para todos los paquetes de su captura. Las seis claves reproducen el digest en el 100\% de los paquetes (179/179 en total), lo que descarta coincidencias. La \autoref{tab:keys} resume cada clave, su índice base-36 y la validación por captura.

\begin{table}[H]
\centering
\footnotesize
\begin{tabular}{l c r l r c}
\hline
\textbf{Captura} & \textbf{Long.} & \textbf{Espacio $36^{n}$} & \textbf{Clave} & \textbf{Índice base-36} & \textbf{Valid.}\\
\hline
06\_01 & 6 & 2{,}176{,}782{,}336     & \texttt{ykusrq}   & 1{,}618{,}816{,}210     & 21/21\\
06\_02 & 6 & 2{,}176{,}782{,}336     & \texttt{miy4or}   & 1{,}673{,}128{,}894     & 57/57\\
07\_01 & 7 & 78{,}364{,}164{,}096    & \texttt{vs7evd8}  & 18{,}253{,}050{,}367    & 27/27\\
07\_02 & 7 & 78{,}364{,}164{,}096    & \texttt{tpx6pje}  & 31{,}666{,}124{,}081    & 24/24\\
07\_03 & 7 & 78{,}364{,}164{,}096    & \texttt{nh3cote}  & 32{,}269{,}346{,}987    & 28/28\\
08\_01 & 8 & 2{,}821{,}109{,}907{,}456 & \texttt{n3s4mmcv} & 2{,}456{,}777{,}905{,}475 & 22/22\\
\hline
\end{tabular}
\caption{Claves recuperadas (todas base-36, keyed-MD5) y validación contra los paquetes autenticados de cada captura.}
\label{tab:keys}
\end{table}

\paragraph{Implicaciones de seguridad.} keyed-MD5 frente a HMAC cambia cómo se calcula un candidato, no cuántos existen. la debilidad es el tamaño del espacio, no la elección del MAC. Para seis caracteres base-36 el espacio es $36^6 = 2{,}176{,}782{,}336$ candidatos, cae en minutos. Un atacante que capture un solo paquete autenticado puede recuperar la clave sin conexión, de modo que el digest ofrece poca protección con claves cortas sobre un alfabeto pequeño. Cada carácter adicional multiplica el espacio por 36, la clave de ocho caracteres demando días de cómputo distribuido (\autoref{sec:bench}).
